% =====================================================================
\chapter{Estadística descriptiva}
\label{tema:1}
% =====================================================================

Antes de construir modelos conviene mirar los datos: resumirlos con unos pocos números y dibujarlos. Este tema presenta las medidas de centralización (media, mediana), de posición (cuantiles y cuartiles), de dispersión (rango, recorrido intercuartílico, varianza y desviación típica), las representaciones gráficas básicas (histograma y diagrama de caja) y las medidas de relación entre dos variables (covarianza y correlación). Se sigue el capítulo~2 de \emph{Introductory Statistics} de OpenStax~\cite{openstax} (§2.2 a §2.7), el capítulo~10 de Rice~\cite{rice}, «Summarizing Data», y, para la covarianza y la correlación, la lectura~7b del curso 18.05 del MIT~\cite{mit1805} y la §7.3 de Blitzstein y Hwang~\cite{bh}.

\section{Datos, población y muestra}

\begin{definicion}[Población, muestra y variable]\label{def:t1-poblacion}
La \emph{población} es el conjunto de todos los individuos u objetos sobre los que se quiere obtener información; una \emph{muestra} es un subconjunto de la población que efectivamente se observa. Una \emph{variable} es una característica que se mide en cada individuo. La variable es \emph{cuantitativa} si toma valores numéricos con los que tiene sentido operar (discreta si toma valores aislados, como un número de errores; continua si puede tomar cualquier valor de un intervalo, como un tiempo) y \emph{cualitativa} o \emph{categórica} si toma valores en un conjunto de categorías.
\end{definicion}

Un conjunto de datos cuantitativos es una lista $x_1,\dots,x_n$ de números reales, donde $n$ es el \emph{tamaño} de la muestra. Los mismos datos ordenados de menor a mayor se escriben $x_{(1)}\le x_{(2)}\le\dots\le x_{(n)}$; $x_{(1)}$ es el mínimo y $x_{(n)}$, el máximo. En este tema los datos son números fijos y las definiciones son puramente descriptivas; a partir del \cref{tema:2} se interpretarán como valores observados de variables aleatorias.

\addtocounter{definicion}{1}%

\section{Medidas de centralización}

\begin{definicion}[Media, mediana y moda]\label{def:t1-centro}
Sean $x_1,\dots,x_n$ datos reales.
\begin{enumerate}
  \item La \emph{media} (aritmética) es $\displaystyle\xbar=\frac1n\sum_{i=1}^{n}x_i$.
  \item La \emph{mediana} es el valor central de los datos ordenados: $x_{((n+1)/2)}$ si $n$ es impar, y $\tfrac12\bigl(x_{(n/2)}+x_{(n/2+1)}\bigr)$ si $n$ es par.
  \item La \emph{moda} es el valor que aparece con mayor frecuencia (puede no ser única).
\end{enumerate}
\end{definicion}

\addtocounter{definicion}{1}%

El ejemplo ilustra que la media es \emph{sensible a los valores atípicos} y la mediana es \emph{robusta}. Cuando la distribución de los datos es asimétrica con cola larga a la derecha (como las longitudes de respuesta, los tiempos o los ingresos), la media suele ser mayor que la mediana~\cite[§2.6]{openstax}. Las dos medidas tienen una caracterización variacional que explica ese comportamiento.

\begin{proposicion}[Propiedades extremales de la media y la mediana]\label{prop:t1-extremal}
Sean $x_1,\dots,x_n$ datos reales.
\begin{enumerate}
  \item $\sum_{i=1}^{n}(x_i-\xbar)=0$.
  \item Para todo $a\in\R$, $\displaystyle\sum_{i=1}^{n}(x_i-a)^{2}=\sum_{i=1}^{n}(x_i-\xbar)^{2}+n(\xbar-a)^{2}$. En consecuencia, la función $a\mapsto\sum_i(x_i-a)^{2}$ alcanza su mínimo únicamente en $a=\xbar$.
  \item La función $a\mapsto\sum_i\abs{x_i-a}$ alcanza su mínimo en la mediana.
\end{enumerate}
\end{proposicion}

\section{Medidas de posición: cuantiles y diagrama de caja}

La mediana divide los datos ordenados en dos mitades. Los cuartiles los dividen en cuatro partes. Hay varias convenciones para calcularlos que dan resultados ligeramente distintos con pocos datos (los programas estadísticos suelen interpolar); en este módulo se usa la siguiente, que es la de~\cite[§2.3]{openstax} y la habitual en el cálculo a mano.

\begin{definicion}[Cuartiles y recorrido intercuartílico]\label{def:t1-cuartiles}
Sean $x_{(1)}\le\dots\le x_{(n)}$ los datos ordenados. La \emph{mitad inferior} está formada por los $\lfloor n/2\rfloor$ primeros datos y la \emph{mitad superior}, por los $\lfloor n/2\rfloor$ últimos (si $n$ es impar, la mediana no pertenece a ninguna de las dos). El \emph{primer cuartil} $Q_1$ es la mediana de la mitad inferior; el \emph{segundo cuartil} $Q_2$ es la mediana; el \emph{tercer cuartil} $Q_3$ es la mediana de la mitad superior. El \emph{recorrido intercuartílico} es $\RIC=Q_3-Q_1$ y el \emph{rango} es $x_{(n)}-x_{(1)}$.
\end{definicion}

Más en general, el \emph{cuantil} de orden $q\in(0,1)$ (o \emph{percentil} $100q$) es un valor que deja a su izquierda aproximadamente una fracción $q$ de los datos; $Q_1$, $Q_2$ y $Q_3$ son los cuantiles de órdenes $0{,}25$, $0{,}5$ y $0{,}75$.

\begin{definicion}[Resumen de cinco números y diagrama de caja]\label{def:t1-caja}
El \emph{resumen de cinco números} es $(x_{(1)},Q_1,Q_2,Q_3,x_{(n)})$. Un dato se considera \emph{atípico} (criterio de Tukey) si está fuera del intervalo
\[
\bigl[\,Q_1-1{,}5\,\RIC,\;Q_3+1{,}5\,\RIC\,\bigr],
\]
cuyos extremos se llaman \emph{vallas}. El \emph{diagrama de caja} dibuja una caja de $Q_1$ a $Q_3$ con una línea en la mediana y dos «bigotes» que llegan hasta el dato más extremo no atípico de cada lado; los atípicos se marcan por separado~\cite[§2.4]{openstax}.
\end{definicion}

\addtocounter{definicion}{1}%

\addtocounter{figure}{1}%

\begin{definicion}[Histograma]\label{def:t1-histograma}
Se divide un intervalo que contiene los datos en \emph{clases} contiguas $[a_0,a_1),[a_1,a_2),\dots,[a_{k-1},a_k]$ y se cuenta la \emph{frecuencia absoluta} $f_j$ de datos en cada clase. El \emph{histograma} dibuja sobre cada clase un rectángulo de altura proporcional a $f_j$ (si las clases tienen igual amplitud) o a $f_j/(a_j-a_{j-1})$ (si no la tienen, de modo que el \emph{área} sea proporcional a la frecuencia). La \emph{frecuencia relativa} es $f_j/n$.
\end{definicion}

El histograma muestra la forma de la distribución (simetría, asimetría, número de modas), mientras que el diagrama de caja resume la posición y la dispersión y señala los atípicos; la \cref{fig:t1-hist-caja} muestra ambos para un mismo conjunto de datos.

\section{Medidas de dispersión}

\begin{definicion}[Varianza y desviación típica muestrales]\label{def:t1-varianza}
La \emph{varianza muestral} de $x_1,\dots,x_n$ ($n\ge 2$) es
\[
s^{2}=\frac{1}{n-1}\sum_{i=1}^{n}(x_i-\xbar)^{2},
\]
y la \emph{desviación típica muestral} es $s=\sqrt{s^{2}}$. La cantidad $\sum_i(x_i-\xbar)^{2}$ se llama \emph{suma de cuadrados} de los datos.
\end{definicion}

\begin{observacion}[¿Por qué $n-1$?]\label{obs:t1-n-1}
El denominador $n-1$ en lugar de $n$ tiene dos justificaciones. La primera es algebraica: las $n$ desviaciones $x_i-\xbar$ suman cero (\cref{prop:t1-extremal}), de modo que solo $n-1$ de ellas pueden elegirse libremente; se dice que la suma de cuadrados tiene $n-1$ \emph{grados de libertad}. La segunda es estadística: si los datos son una muestra aleatoria de una población con varianza $\sigma^{2}$, la versión con $n-1$ acierta «en promedio», es decir, es un estimador insesgado de $\sigma^{2}$, mientras que la versión con $n$ la subestima sistemáticamente. Esto se demuestra en el \cref{tema:2} (\cref{teo:t2-s2}). Con denominador $n$ se obtiene la \emph{varianza poblacional} de los propios datos, $\frac1n\sum_i(x_i-\xbar)^{2}$, que es la que se usa cuando los datos \emph{son} toda la población o cuando se normaliza un vector, como en la normalización por capas de una red (módulo~11).
\end{observacion}

\begin{proposicion}[Fórmula de cálculo]\label{prop:t1-abreviada}
$\displaystyle\sum_{i=1}^{n}(x_i-\xbar)^{2}=\sum_{i=1}^{n}x_i^{2}-n\xbar^{2}$, y por tanto $s^{2}=\dfrac{\sum_i x_i^{2}-n\xbar^{2}}{n-1}$.
\end{proposicion}

\addtocounter{definicion}{1}%

\begin{proposicion}[Cambios de escala y de origen]\label{prop:t1-lineal}
Si $y_i=ax_i+b$ con $a,b\in\R$, entonces $\ybar=a\xbar+b$, $s_y^{2}=a^{2}s_x^{2}$ y $s_y=\abs{a}\,s_x$. Si $a>0$, la mediana y los cuartiles se transforman igual que la media y $\RIC_y=a\,\RIC_x$.
\end{proposicion}

\begin{definicion}[Tipificación]\label{def:t1-z}
Si $s>0$, la \emph{puntuación tipificada} (o puntuación $z$) del dato $x_i$ es $z_i=(x_i-\xbar)/s$. Por la \cref{prop:t1-lineal}, los datos tipificados tienen media $0$ y desviación típica $1$; $z_i$ mide a cuántas desviaciones típicas está $x_i$ de la media.
\end{definicion}

\section{Datos bivariantes: covarianza y correlación}

Cuando en cada individuo se miden dos variables se obtienen pares $(x_1,y_1),\dots,(x_n,y_n)$, que se representan en un \emph{diagrama de dispersión}.

\begin{definicion}[Covarianza y correlación muestrales]\label{def:t1-cov}
La \emph{covarianza muestral} y el \emph{coeficiente de correlación} (de Pearson) son
\[
s_{xy}=\frac{1}{n-1}\sum_{i=1}^{n}(x_i-\xbar)(y_i-\ybar),
\qquad
r=\frac{s_{xy}}{s_x\,s_y}=\frac{\sum_i(x_i-\xbar)(y_i-\ybar)}{\sqrt{\sum_i(x_i-\xbar)^{2}}\,\sqrt{\sum_i(y_i-\ybar)^{2}}},
\]
donde se supone $s_x,s_y>0$.
\end{definicion}

El signo de cada producto $(x_i-\xbar)(y_i-\ybar)$ es positivo cuando el punto está en el cuadrante superior derecho o inferior izquierdo respecto de $(\xbar,\ybar)$ y negativo en los otros dos; $s_{xy}>0$ indica que predominan los primeros, es decir, una tendencia creciente. La covarianza depende de las unidades; la correlación no (\cref{eje:t1-r-invariante}).

\begin{teorema}[Acotación de la correlación]\label{teo:t1-r}
Se cumple $-1\le r\le 1$. Además, $\abs{r}=1$ si y solo si existen $a\neq 0$ y $b$ tales que $y_i=ax_i+b$ para todo $i$, y entonces el signo de $r$ es el de $a$.
\end{teorema}

\addtocounter{definicion}{1}%

Los mismos conceptos existen para variables aleatorias. Si $(X,Y)$ es un vector aleatorio con varianzas finitas y positivas, su \emph{covarianza} y su \emph{correlación} son
\[
\Cov(X,Y)=\E\bigl[(X-\E X)(Y-\E Y)\bigr]=\E[XY]-\E[X]\E[Y],
\qquad
\Cor(X,Y)=\frac{\Cov(X,Y)}{\sqrt{\Var X\,\Var Y}},
\]
y también $\abs{\Cor(X,Y)}\le1$~\cite[§7.3]{bh}. Las fórmulas muestrales de la \cref{def:t1-cov} son las análogas con los datos (salvo el factor $n/(n-1)$, que se cancela en $r$).

\addtocounter{definicion}{1}%

\addtocounter{definicion}{1}%

